% Lösungen Maßstab, Ähnlichkeit und Strahlensätze (zusammenbau v0.4, Heft)
\einheitenkopf{Maßstab, Ähnlichkeit und Strahlensätze}

\erg{58}{a) kleiner \quad b) $20\,\text{cm}$ \quad c) $6\,\text{cm}$}
\erg{59}{a) $28 : 40 = 0{,}7\,\text{m} = 70\,\text{cm}$; $5{,}6 : 40 = 0{,}14\,\text{m} = 14\,\text{cm}$ \quad b) $36 : 80 = 0{,}45\,\text{m} = 45\,\text{cm}$; $9{,}6 : 80 = 0{,}12\,\text{m} = 12\,\text{cm}$ \quad c) $0{,}92 \cdot 25 = 23\,\text{m}$; $0{,}64 \cdot 25 = 16\,\text{m}$ \quad d) $0{,}45 \cdot 30 = 13{,}5\,\text{m}$; $0{,}62 \cdot 30 = 18{,}6\,\text{m}$}
\erg{60}{a) $1 : 40$ \quad b) $7$ Kästchen \quad c) Breite $4 \cdot 2 = 8$ Kästchen, Höhe $3 \cdot 2 = 6$ Kästchen}

60c)

\rechteck{8}{6}

\erg{61}{a) Draufsicht: Kreis mit $d = 66\,\text{cm}$ im Rechteck $60\,\text{cm}$ mal $85\,\text{cm}$; z. B. Maßstab $1 : 10$ (Rechteck $6\,\text{cm}$ mal $8{,}5\,\text{cm}$, Kreis $d = 6{,}6\,\text{cm}$); nein, der Durchmesser $66\,\text{cm}$ ist größer als die Plattenbreite $60\,\text{cm}$ \quad b) Draufsicht: Kreis mit $d = 3\,\text{m}$ im Rechteck $2{,}8\,\text{m}$ mal $4\,\text{m}$; z. B. Maßstab $1 : 40$ (Rechteck $7\,\text{cm}$ mal $10\,\text{cm}$, Kreis $d = 7{,}5\,\text{cm}$); nein, der Durchmesser $3\,\text{m}$ ist größer als die Planenbreite $2{,}8\,\text{m}$}
